\section{Conclusions}

We have developed a hybrid deterministic-stochastic model to simulate collapse-recovery dynamics across ten plausible futures for Earth-originating civilization, and used these simulations to quantify the duty cycles of technological civilizations. Our principal findings are as follows.

First, collapse-recovery dynamics produce a wide range of duty cycles across plausible civilizational futures, with $D_c$ spanning from approximately 0.38 to 1.00 across the ten scenarios, with outcomes shaped by the interplay of governance structure, resource pressure, and hazard exposure.

Second, the duty-cycle framework reveals that several model parameters correspond to actionable levers for civilizational resilience (Section~\ref{earth-implications}). Our sensitivity analysis demonstrates that $r_f$ and $\delta$ consistently emerge as the most impactful parameters across all collapse-prone scenarios, and that modest shifts in these levers can qualitatively alter a civilization's long-term trajectory, suggesting that targeted resilience investments could yield disproportionate returns in civilizational continuity.

Third, the effective detectability duration $L_{\mathrm{eff}} = D_c \, T_{\mathrm{span}}$ provides a more realistic replacement for the Drake Equation's static longevity term $L$. Combined with power-law longevity distributions that disfavor long-lived technosignatures \citep{Balbi2024-bg}, low duty cycles imply a doubly diminished probability of detection, helping explain the observed absence of extraterrestrial signals.

Several avenues for future work emerge from this study. Incorporating exergy constraints and non-linear growth functions informed by thermodynamic limits to growth \citep{Haqq-Misra2025-lum} could yield more physically grounded simulations. Introducing feedback between technological capacity and accessible resource stocks would address a key simplification of the current model. Finally, extending the framework to account for per-technosignature observation probabilities and for the possibility that autonomous remnant technologies may cycle through their own active and dormant states would add further realism to the detectability analysis.

More broadly, our results help reconcile seemingly contradictory perspectives on civilizational longevity. One view posits that technological civilizations are fundamentally fragile and prone to self-destruction soon after emergence \citep{Vinn2024-pi}. Another suggests that once civilizations surpass key thresholds of resilience and coordination, they may persist indefinitely, steadily accumulating across cosmic time \citep{Grinspoon2009-cg}. Both patterns emerge naturally in our simulations---not from differences in exogenous shocks, but from internal sociotechnical structure. The long-term fate of a civilization, it appears, is less a matter of luck than of design.